\chapter{Case Study: Building an FST-Based Spellchecker}
\label{ch:fst-spellchecker}

A useful spellchecker must distinguish an absent correction from an incomplete search.
It must also avoid returning the same word several times merely because several edit scripts produce it.
The case study builds a bounded, executable FST search with deterministic ranking and an independent distance oracle.
The implementation is \code{strings.py} in \path{example-code/string-labs/}.
It requires only Python's standard library.

\section{Specify what a correction means}
\label{sec:ch80-contract}
The dictionary consists of nonempty strings with positive integer frequencies.
Both dictionary words and queries use the NFC--casefold--NFC key policy from Chapter~\ref{ch:strings-representation}.
Distance counts insertion, deletion, and substitution of normalized Unicode code points, not bytes or grapheme clusters.
Adjacent swaps cost two when they can be expressed by two ordinary edits.
The search radius ranges from zero through three, and query and dictionary word lengths are checked against 128 code points before and after normalization.

These choices make the algorithm reproducible but do not make it a complete natural-language spelling system.
There is no context model, language identification, tokenization of documents, morphology generator, keyboard adjacency model, or locale-specific casing.
The tiny checked-in dictionary is an original demonstration, not a benchmark corpus.
The public result contains canonical word, minimum edit cost, and accumulated dictionary frequency.

Table~\ref{tab:ch80-ranking} works query \code{ct} for four dictionary words.
At radius one the complete candidates are cat, cot, and cut, ordered by decreasing frequency because their distances tie.
Coat is excluded at radius one despite its frequency exceeding cot's.
Edit distance is the primary key, frequency the secondary key, and canonical lexical spelling the final deterministic tie breaker.

\begin{table}[htbp]
\caption{Frequency ranks equal-distance candidates but cannot admit an out-of-radius word.}
\label{tab:ch80-ranking}
\centering
\begin{tabular}{lrrl}
\toprule
Word & Distance from ct & Frequency & Radius-one outcome \\
\midrule
cat & 1 & 30 & first \\
coat & 2 & 20 & excluded \\
cot & 1 & 10 & second \\
cut & 1 & 5 & third \\
\bottomrule
\end{tabular}
\end{table}

\section{Compile the dictionary identity transducer}
\label{sec:ch80-build}
Each trie edge labeled by character $c$ is interpreted as the identity-transducer arc $c:c$.
Final states record word frequency.
Thus the implementation genuinely composes an edit relation with a finite-state dictionary relation; it does not generate a set of misspellings and repeatedly perform independent dictionary lookups.
The composition is specialized and implicit rather than delegated to a general-purpose FST package.

Algorithm~\ref{alg:ch80-build} first merges equal canonical keys, then inserts them into the trie.
Frequency merging is addition: entries \code{Straße} with frequency three and \code{STRASSE} with frequency four create one final key \code{strasse} with frequency seven.
Original display variants are not retained by this lab.
A production UI could add a separate preferred-spelling field without changing distance computation.

\begin{algorithm}[htbp]
\caption{Compile bounded normalized dictionary entries into an identity transducer.}
\label{alg:ch80-build}
\begin{pseudocode}
\Require finite word/frequency entries; positive frequencies; length policy
\Ensure trie with one final entry per canonical spelling
\ForAll{entries}
\State validate type, frequency, and original extent; construct canonical key
\State validate canonical extent; add frequency to that key's total
\EndFor
\State allocate root state
\ForAll{canonical words in lexical order}
\State follow or create one identity arc per symbol
\State mark final state and store merged frequency
\EndFor
\end{pseudocode}
\end{algorithm}

Ignoring normalization and sorting, transition insertion costs expected $O(L)$ for $L$ total canonical symbols with hash-map children.
Sorting $W$ keys costs comparisons whose cost includes their common prefixes.
The implementation also stores each state's full prefix string for simple output reconstruction.
That extra storage can be $O(LM)$ with maximum word length $M$, rather than the $O(L)$ edge storage of a parent-pointer trie.
Here $M\le128$ is a fixed contract bound; a general dictionary library should use parent links or compressed output storage instead of hiding that factor.

Figure~\ref{fig:ch80-components} separates dictionary construction from query-time composition.
No transducer metadata is inferred from a spelling query.
The same immutable trie serves multiple searches, each with its own queue and distances.
The example does not mutate the dictionary concurrently with a search.

\begin{figure}[htbp]
\centering
\begin{tikzpicture}[node distance=12mm]
\node[basic block] (data) {UTF-8 TSV dictionary};
\node[basic block,below=of data] (trie) {normalize, merge, build trie FST};
\node[basic block,below=of trie] (search) {query-specific implicit edit composition};
\node[basic block,below=of search] (rank) {minimum costs, then frequency ranking};
\draw[cfg edge] (data)--(trie);\draw[cfg edge] (trie)--(search);\draw[cfg edge] (search)--(rank);
\end{tikzpicture}
\caption{A persistent dictionary transducer is combined with fresh per-query search state.}
\label{fig:ch80-components}
\end{figure}

\section{Keep best states without losing output words}
\label{sec:ch80-best}
The priority queue stores triples consisting of cost, input position, and trie state ID.
A best-distance map rejects a new path unless it improves the cost at that exact product state.
When a stale queue entry is popped, it is ignored.
Because each trie state has one root path, all paths to $(i,q)$ have emitted the same candidate prefix.
Their remaining legal continuations are identical, so a more expensive path cannot produce a better final candidate.

Insertion advances trie depth, deletion advances input position, and copy/substitution advances both.
Consequently the specialized product has no cycle even though the abstract edit transducer can loop.
Nonnegative weights also justify Dijkstra's settlement order.
At a final state with all input consumed, the popped cost is the minimum distance for that dictionary word.
The independent dynamic-programming oracle checks precisely this semantic result.

It is tempting to stop after the first $K$ final states have been popped.
That is unsafe for the chosen ranking: a later state can have the same distance but a higher frequency.
The implementation explores the complete within-radius reachable product, sorts accepted words by the full ranking key, and then truncates to the requested result count.
An optimized early-stop variant would need a lower bound on the full remaining ranking order, not merely edit distance.

\section{Make resource exhaustion observable}
\label{sec:ch80-budget}
Each nonstale product-state expansion charges the configured budget.
If the next expansion would exceed it, search raises \code{BudgetExceeded} rather than returning a partial list labeled as complete.
The CLI writes no success JSON in that case and exits with status two.
Algorithm~\ref{alg:ch80-query} describes this publication boundary.

\begin{algorithm}[htbp]
\caption{Publish spellchecking results only after complete bounded search.}
\label{alg:ch80-query}
\begin{pseudocode}
\Require dictionary FST, query, radius, positive result limit and expansion budget
\Ensure ranked complete top results, or explicit failure without success output
\State validate and canonicalize query
\State run implicit product search, retaining all within-radius accepted words
\If{state expansion budget is exhausted}
\State report incomplete search; do not publish candidate JSON
\EndIf
\State sort by cost, negative frequency, and canonical spelling
\State publish the requested prefix of the sorted candidate list
\end{pseudocode}
\end{algorithm}

The expansion budget is not a wall-clock or total-allocation bound.
A state with many outgoing dictionary edges can generate many transitions, and dictionary construction itself has separate cost.
The word-length bound controls path depth but not vocabulary size.
An exposed network service would additionally need dictionary-size limits, request-size limits before decoding, time accounting, and cancellation that does not corrupt shared dictionary state.

Table~\ref{tab:ch80-failures} distinguishes valid empty results from failures.
A query containing no nearby word returns an empty JSON array after complete search.
An invalid radius, malformed UTF-8 dictionary, or invalid frequency is a rejected request or dictionary, not evidence that no correction exists.

\begin{table}[htbp]
\caption{Absence of a candidate and absence of a complete search are different outcomes.}
\label{tab:ch80-failures}
\centering
\begin{tabular}{ll}
\toprule
Condition & Public outcome \\
\midrule
complete search; no word in radius & empty successful result \\
complete search; several edit paths & one minimum-cost result per word \\
expansion budget exceeded & explicit incomplete-search error \\
invalid radius or frequency & validation error \\
normalized key collision & merged frequency, one canonical entry \\
\bottomrule
\end{tabular}
\end{table}

\section{Verify against an independent formulation}
\label{sec:ch80-tests}
The lab builds every nonempty binary word of length at most four as a dictionary: 30 words.
It queries every ternary word of length at most four, including empty: 121 queries.
For radii zero through three, it compares the entire ranked result against row-based Levenshtein distance computed independently for each dictionary word.
There are 484 complete candidate-set comparisons and 14,520 query/radius/dictionary-word obligations.
This tests the epsilon insertion and deletion directions, duplicate-path merging, radius pruning, and final acceptance condition together.

Additional tests cover canonical equivalence, case-fold expansion, frequency merging, the worked ranking, invalid bounds, and CLI success/failure statuses.
The exact matcher is checked separately against direct search rather than reused as its own oracle.
Figure~\ref{fig:ch80-oracle} shows the independent paths.
Neither finite testing nor agreement with row DP establishes linguistic usefulness; that would require a separately curated evaluation set and an explicitly chosen quality metric.

\begin{figure}[htbp]
\centering
\begin{tikzpicture}[node distance=13mm and 12mm]
\node[basic block] (input) {query and finite dictionary};
\node[basic block,below left=of input] (fst) {FST product search};
\node[basic block,below right=of input] (dp) {per-word row DP};
\node[basic block,below=of input,yshift=-18mm] (cmp) {compare complete ranked sets};
\draw[cfg edge] (input)--(fst);\draw[cfg edge] (input)--(dp);
\draw[cfg edge] (fst)--(cmp);\draw[cfg edge] (dp)--(cmp);
\end{tikzpicture}
\caption{The oracle computes distances without using the product-state search implementation.}
\label{fig:ch80-oracle}
\end{figure}

\section{Exercises}
\label{sec:ch80-exercises}
\begin{enumerate}
\item\label{ex:ch80-ranking} Rank the four words in Table~\ref{tab:ch80-ranking} at radius two and limit three.
\item\label{ex:ch80-collision} Merge dictionary entries \code{Straße}:3 and \code{STRASSE}:4; give the radius-zero result for query \code{strasse}.
\item\label{ex:ch80-stop} Give a reason that stopping at the first final state can violate frequency tie-breaking even though the state's distance is minimal.
\item\label{ex:ch80-incomplete} Explain why returning an empty list after exhausting a budget of one expansion would violate the documented contract for query \code{cat} in a dictionary containing cat.
\end{enumerate}

\section{Further reading}
\label{sec:ch80-reading}
Mohri's weighted-automata treatment supplies the general framework behind the specialized composition~\cite{mohri-2009-weighted}.
The lab README records executable commands, limits, and the independent tests; its solutions file answers the exercises from all four string chapters.
